\clearpage
\begin{trapbox}
\textbf{\textcolor{CTrap}{易错一｜看到顶点值就直接当成区间最值}}
\begin{center}\begin{tikzpicture}[font=\small,>=stealth,line join=round]
\draw[gray!40,rounded corners=3pt] (0,0) rectangle (14.2,8.3);\node at (7.1,7.700000000000001) {顶点虽低，却不在允许的区间中};
\begin{scope}[shift={(1.8777777777777778,2.2249999999999996)},x=1.111111111111111cm,y=0.575cm]\draw[->,gray] (-0.7,0)--(4.7,0) node[right]{$x$};\draw[->,gray] (0,-1)--(0,7) node[above]{$y$};\draw[gray,thick,solid,domain=-0.2:4.2,samples=180] plot(\x,{(\x-2)^2+1});\draw[blue,thick,very thick,domain=3:4,samples=180] plot(\x,{(\x-2)^2+1});\fill[orange!85!black] (3,2) circle(2.2pt);\fill[orange!85!black] (4,5) circle(2.2pt);\draw[violet,dashed] (2,1.05)--(2,6.5);\draw[orange!85!black,very thick] (3,-.5)--(4,-.5);\node[below] at (2,1){$(2,1)$};\node[right] at (3,2){$(3,2)$};\node[above left=3pt] at (4,5){$(4,5)$};\foreach \v in {2,3,4}{\node[below=5pt] at (\v,-.5){$\v$};}\end{scope}\node[anchor=west,align=left,text width=5.0cm] at (8.4,4.0) {$f=(x-2)^2+1$\\[9pt]允许 $x\in[3,4]$\\[9pt]$2\notin[3,4]$\\[9pt]最近 $x=3$：最小值2\\[9pt]最远 $x=4$：最大值5};
\node[text width=12.6cm,align=center] at (7.1,.65) {灰线仅供定位；蓝色粗线是允许段，不能到灰色顶点处取值。};
\end{tikzpicture}\end{center}



\textbf{错误想法：} $f(x)=(x-2)^2+1$ 的最小值是 $1$，所以在 $[3,4]$ 上也最小为 $1$。

\textbf{为什么错：} 取到 $1$ 必须让 $x=2$，但 $2\notin[3,4]$。

\textbf{正确做法：} 只在允许范围内找，最小值是 $f(3)=2$。

\medskip

\end{trapbox}
\clearpage
\begin{trapbox}
\textbf{\textcolor{CTrap}{易错二｜开口向下还套“远的更大”}}
\begin{center}\begin{tikzpicture}[font=\small,>=stealth,line join=round]
\draw[gray!40,rounded corners=3pt] (0,0) rectangle (14.2,8.3);\node at (7.1,7.700000000000001) {例3：开口向下，最远的右端取得最小值};
\begin{scope}[shift={(2.0545454545454547,2.8)},x=1.3636363636363635cm,y=0.575cm]\draw[->,gray] (-0.7,0)--(3.7,0) node[right]{$x$};\draw[->,gray] (0,-2)--(0,6) node[above]{$y$};\draw[blue,thick,very thick,domain=0:3,samples=180] plot(\x,{-(\x-1)^2+4});\fill[orange!85!black] (0,3) circle(2.2pt);\fill[orange!85!black] (1,4) circle(2.2pt);\fill[orange!85!black] (3,0) circle(2.2pt);\node[left] at (0,3){$(0,3)$};\node[above] at (1,4){$(1,4)$};\node[below] at (3,0){$(3,0)$};\end{scope}\node[anchor=west,align=left,text width=5.0cm] at (8.4,4.0) {$f=-(x-1)^2+4$\\[9pt]$x\in[0,3]$\\[9pt]轴 $x=1$\\[9pt]最近取最大值4\\[9pt]最远$x=3$，最小值0};
\node[text width=12.6cm,align=center] at (7.1,.65) {左端距离1、右端距离2；更远的端点反而更低。};
\end{tikzpicture}\end{center}



\textbf{反例：} $f(x)=-(x-1)^2+4$，在 $[0,3]$ 上，离 $1$ 最远的端点 $3$ 取得的是\textbf{最小值 $0$}，不是最大值。

\textbf{正确做法：} 先看 $a$ 的正负：$a>0$ 近小远大；$a<0$ 近大远小。

\medskip

\end{trapbox}
\clearpage
\begin{trapbox}
\textbf{\textcolor{CTrap}{易错三｜顶点在区间内，却以为两端点同高}}
\begin{center}\begin{tikzpicture}[font=\small,>=stealth,line join=round]
\draw[gray!40,rounded corners=3pt] (0,0) rectangle (14.2,8.3);\node at (7.1,7.700000000000001) {顶点在区间内，并不代表两端点同高};
\begin{scope}[shift={(1.8777777777777778,2.307142857142857)},x=1.111111111111111cm,y=0.32857142857142857cm]\draw[->,gray] (-0.7,0)--(4.7,0) node[right]{$x$};\draw[->,gray] (0,-2)--(0,12) node[above]{$y$};\draw[blue,thick,very thick,domain=0:4,samples=180] plot(\x,{(\x-1)^2});\fill[orange!85!black] (0,1) circle(2.2pt);\fill[orange!85!black] (1,0) circle(2.2pt);\fill[orange!85!black] (4,9) circle(2.2pt);\draw[violet,dashed] (1,.1)--(1,11);\node[left] at (0,1){$(0,1)$};\node[above] at (4,9){$(4,9)$};\node[below] at (1,0){$1$};\end{scope}\node[anchor=west,align=left,text width=5.0cm] at (8.4,4.0) {$f=(x-1)^2$\\[9pt]$x\in[0,4]$\\[9pt]$h=1$在区间内\\[9pt]却不在中点2\\[9pt]$f(0)=1\ne f(4)=9$};
\node[text width=12.6cm,align=center] at (7.1,.65) {只有轴在区间中点，两端到轴的距离才相等，函数值才相等。};
\end{tikzpicture}\end{center}



\textbf{反例：} $f(x)=(x-1)^2$，$x\in[0,4]$，对称轴在区间内，但 $f(0)=1$、$f(4)=9$ 并不相等。

\textbf{正确做法：} 只有对称轴\textbf{正好在区间中点}时，两端才同高：
\[
 f(m)=f(n)\Longleftrightarrow h=\frac{m+n}{2}\quad(a\ne0,\ m<n).
\]

\medskip

\end{trapbox}
\clearpage
\begin{trapbox}
\textbf{\textcolor{CTrap}{易错四｜以为求最值一定得算出三个候选点}}
\begin{center}\begin{tikzpicture}[font=\small,>=stealth,line join=round]
\draw[gray!40,rounded corners=3pt] (0,0) rectangle (14.2,8.6);\node at (7.1,8.0) {比远近：看对称轴在区间中点的哪一侧};
\begin{scope}[shift={(2,6.3)},x=1.1cm]\draw[->,gray] (-.5,0)--(8.8,0);\draw[blue,thick] (0,0)--(8,0);\foreach \v in {0,8}{\fill[blue] (\v,0) circle(2pt);\node[below=4pt] at (\v,0){$\v$};}\draw[violet,dashed] (4,-.2)--(4,.15);\fill[orange!85!black] (3,0) circle(2.5pt);\node[above=6pt] at (3,0){$h=3$};\node[anchor=west] at (0,1){左端距离3，右端距离5：右端更远};\end{scope}\begin{scope}[shift={(2,4.4)},x=1.1cm]\draw[->,gray] (-.5,0)--(8.8,0);\draw[blue,thick] (0,0)--(8,0);\foreach \v in {0,8}{\fill[blue] (\v,0) circle(2pt);\node[below=4pt] at (\v,0){$\v$};}\draw[violet,dashed] (4,-.2)--(4,.15);\fill[orange!85!black] (6,0) circle(2.5pt);\node[above=6pt] at (6,0){$h=6$};\node[anchor=west] at (0,1){左端距离6，右端距离2：左端更远};\end{scope}\begin{scope}[shift={(2,2.5)},x=1.1cm]\draw[->,gray] (-.5,0)--(8.8,0);\draw[blue,thick] (0,0)--(8,0);\foreach \v in {0,8}{\fill[blue] (\v,0) circle(2pt);\node[below=4pt] at (\v,0){$\v$};}\draw[violet,dashed] (4,-.2)--(4,.15);\fill[orange!85!black] (4,0) circle(2.5pt);\node[above=6pt] at (4,0){$h=4$};\node[anchor=west] at (0,1){左端距离4，右端距离4：两端同远};\end{scope}
\node[text width=12.6cm,align=center] at (7.1,.65) {同尺度；紫线为中点4，轴位于中点时两端才同远。};
\end{tikzpicture}\end{center}



\textbf{问题：} 许多题先算 $f(m),f(n),f(h)$，既费时间又可能把区间外的 $h$ 算进去。

\textbf{更快：} 先画数轴判断最近、最远；只代入\textbf{真正决定最值的两个位置}。区间中点还可以直接判断哪个端点更远。

\medskip

\end{trapbox}
\clearpage
\begin{trapbox}
\textbf{\textcolor{CTrap}{易错五｜忘记区间是否包含端点}}
\begin{center}\begin{tikzpicture}[font=\small,>=stealth,line join=round]
\draw[gray!40,rounded corners=3pt] (0,0) rectangle (14.2,8.8);\node at (7.1,8.200000000000001) {易错五：边界能接近，不等于能取到};
\node at (3.1500000000000004,7.3) {$x\in[0,1]$};\begin{scope}[shift={(1.61875,3.1333333333333333)},x=3.0625cm,y=2.111111111111111cm]\draw[->,gray] (-0.3,0)--(1.3,0) node[right]{$x$};\draw[->,gray] (0,-0.3)--(0,1.5) node[above]{$y$};\draw[blue,thick,solid,domain=0:1,samples=180] plot(\x,{(\x)^2});\fill[orange!85!black] (0,0) circle(2pt);\fill[orange!85!black] (1,1) circle(2pt);\end{scope}\node[align=center,text width=5.8cm] at (3.1500000000000004,1.7) {最小0，最大1\\两端都能取到};\node at (9.95,7.3) {$x\in[0,1)$};\begin{scope}[shift={(8.41875,3.1333333333333333)},x=3.0625cm,y=2.111111111111111cm]\draw[->,gray] (-0.3,0)--(1.3,0) node[right]{$x$};\draw[->,gray] (0,-0.3)--(0,1.5) node[above]{$y$};\draw[teal,thick,solid,domain=0:1,samples=180] plot(\x,{(\x)^2});\fill[orange!85!black] (0,0) circle(2pt);\draw[teal,fill=white,thick] (1,1) circle(3pt);\end{scope}\node[align=center,text width=5.8cm] at (9.95,1.7) {最小0，无最大值\\高度1只能接近};
\node[text width=12.6cm,align=center] at (7.1,.65) {两图同尺度，均为$y=x^2$；实心点可取，空心点不可取。};
\end{tikzpicture}\end{center}



\textbf{对比：}
\begin{itemize}[leftmargin=*,itemsep=0.2em]
\item $f(x)=x^2$，$x\in[0,1]$：最小值 $0$，最大值 $1$，都能取到。
\item $f(x)=x^2$，$x\in[0,1)$：最小值仍为 $0$；$x$ 能越来越靠近 $1$，但不能等于 $1$，因此\textbf{没有最大值}。
\end{itemize}
\textbf{正确做法：} 本结论针对\textbf{闭区间}。看见 $($ 或 $)$，要重新确认边界值到底能不能取到。

\medskip

\end{trapbox}
\clearpage
\begin{trapbox}
\textbf{\textcolor{CTrap}{易错六｜把 $h=m$ 或 $h=n$ 当成“顶点在区间外”}}
\begin{center}\begin{tikzpicture}[font=\small,>=stealth,line join=round]
\draw[gray!40,rounded corners=3pt] (0,0) rectangle (14.2,8.8);\node at (7.1,8.200000000000001) {易错六：轴恰在端点，仍能取到顶点值};
\node at (3.1500000000000004,7.3) {$f=(x-1)^2+2$};\begin{scope}[shift={(1.3125,2.723529411764706)},x=1.225cm,y=0.44705882352941173cm]\draw[->,gray] (-0.5,0)--(3.5,0) node[right]{$x$};\draw[->,gray] (0,-0.5)--(0,8) node[above]{$y$};\draw[blue,thick,solid,domain=1:3,samples=180] plot(\x,{(\x-1)^2+2});\fill[orange!85!black] (1,2) circle(2pt);\fill[orange!85!black] (3,6) circle(2pt);\end{scope}\node[align=center,text width=5.8cm] at (3.1500000000000004,1.7) {$[1,3]$，轴等于左端\\最小2，在$x=1$};\node at (9.95,7.3) {$f=(x-1)^2+2$};\begin{scope}[shift={(9.3375,2.723529411764706)},x=1.225cm,y=0.44705882352941173cm]\draw[->,gray] (-1.5,0)--(2.5,0) node[right]{$x$};\draw[->,gray] (0,-0.5)--(0,8) node[above]{$y$};\draw[teal,thick,solid,domain=-1:1,samples=180] plot(\x,{(\x-1)^2+2});\fill[orange!85!black] (-1,6) circle(2pt);\fill[orange!85!black] (1,2) circle(2pt);\end{scope}\node[align=center,text width=5.8cm] at (9.95,1.7) {$[-1,1]$，轴等于右端\\最小2，在$x=1$};
\node[text width=12.6cm,align=center] at (7.1,.65) {两图同尺度；闭区间的端点包含在内，判断式中的等号不能漏。};
\end{tikzpicture}\end{center}



\textbf{正确做法：} 闭区间包含端点，$m\le h\le n$ 包括等号。顶点恰好落在端点时，顶点值仍可以取到。

\medskip

\end{trapbox}
\clearpage
\begin{trapbox}
\textbf{\textcolor{CTrap}{易错七｜含参数时没有检查区间是否合法}}
\begin{center}\begin{tikzpicture}[font=\small,>=stealth,line join=round]
\draw[gray!40,rounded corners=3pt] (0,0) rectangle (14.2,8.3);\node at (7.1,7.700000000000001) {易错七：先检查左端是否小于右端};
\begin{scope}[shift={(3,5.5)},x=1.5cm]\draw[->,gray] (-1.5,0)--(4,0);\draw[blue,very thick] (0,0)--(3,0);\foreach \v/\L in {0/0,3/m}{\fill[blue] (\v,0) circle(2pt);\node[below=5pt] at (\v,0){$\L$};}\node[above=10pt] at (1.5,0){$m>0$：左端0小于右端m};\end{scope}\begin{scope}[shift={(3,3)},x=1.5cm]\draw[->,gray] (-1.5,0)--(4,0);\fill[orange!85!black] (-1,0) circle(2pt);\fill[blue] (0,0) circle(2pt);\node[below=5pt] at (-1,0){$m<0$};\node[below=5pt] at (0,0){$0$};\node[anchor=west] at (.5,.55){此时m在0左侧};\end{scope}
\node[text width=12.6cm,align=center] at (7.1,.65) {$[0,m]$若要求两端不同，就应有$m>0$；$m=0$时只剩一个点。};
\end{tikzpicture}\end{center}



\textbf{例：} 写成 $[0,m]$，通常需要 $m\ge0$；如果题目明确是两个不同端点构成的区间，就要 $m>0$。不能一边说 $m<0$，一边仍把 $[0,m]$ 当成从左到右的普通区间来使用。

\textbf{正确做法：} 先确定参数允许的范围，再谈顶点、区间中点和两端点。

\medskip

\end{trapbox}
\clearpage
\begin{trapbox}
\textbf{\textcolor{CTrap}{易错八｜看错“远近”与“左右”}}
\begin{center}\begin{tikzpicture}[font=\small,>=stealth,line join=round]
\draw[gray!40,rounded corners=3pt] (0,0) rectangle (14.2,8.6);\node at (7.1,8.0) {比远近：看对称轴在区间中点的哪一侧};
\begin{scope}[shift={(2,6.3)},x=1.1cm]\draw[->,gray] (-.5,0)--(8.8,0);\draw[blue,thick] (0,0)--(8,0);\foreach \v in {0,8}{\fill[blue] (\v,0) circle(2pt);\node[below=4pt] at (\v,0){$\v$};}\draw[violet,dashed] (4,-.2)--(4,.15);\fill[orange!85!black] (3,0) circle(2.5pt);\node[above=6pt] at (3,0){$h=3$};\node[anchor=west] at (0,1){左端距离3，右端距离5：右端更远};\end{scope}\begin{scope}[shift={(2,4.4)},x=1.1cm]\draw[->,gray] (-.5,0)--(8.8,0);\draw[blue,thick] (0,0)--(8,0);\foreach \v in {0,8}{\fill[blue] (\v,0) circle(2pt);\node[below=4pt] at (\v,0){$\v$};}\draw[violet,dashed] (4,-.2)--(4,.15);\fill[orange!85!black] (6,0) circle(2.5pt);\node[above=6pt] at (6,0){$h=6$};\node[anchor=west] at (0,1){左端距离6，右端距离2：左端更远};\end{scope}\begin{scope}[shift={(2,2.5)},x=1.1cm]\draw[->,gray] (-.5,0)--(8.8,0);\draw[blue,thick] (0,0)--(8,0);\foreach \v in {0,8}{\fill[blue] (\v,0) circle(2pt);\node[below=4pt] at (\v,0){$\v$};}\draw[violet,dashed] (4,-.2)--(4,.15);\fill[orange!85!black] (4,0) circle(2.5pt);\node[above=6pt] at (4,0){$h=4$};\node[anchor=west] at (0,1){左端距离4，右端距离4：两端同远};\end{scope}
\node[text width=12.6cm,align=center] at (7.1,.65) {同尺度；紫线为中点4，轴位于中点时两端才同远。};
\end{tikzpicture}\end{center}



\textbf{错误想法：} 右端点一定更远，或者图像最右边一定最高。

\textbf{正确做法：} 离对称轴更远，不一定更靠右。例如 $h=6$、区间 $[0,8]$，左端点距离为 $6$，右端点距离为 $2$，左端才更远。
\end{trapbox}
